% GENERATED FILE. Edit canonical lessons, then run npm run generate:books.
% canonical-source-hash: fnv1a64:0ab6c7821e5a9f02
% canonical-lessons: JA-C72-gogo, JA-C72-yonaka, JA-C72-asagata, JA-C72-yuuhan, JA-C72-asahi

\chapter{Afternoon, Middle of the night, Early morning, Evening meal, Morning sun}
\label{ch:ja-afternoon-middle-of-the-night-early-morning-evening-meal-morning-sun}
\hlchaptermodality{72}

\begin{chapteropening}
\textbf{By the end of this chapter:} \emph{I can say the afternoon, the middle of the night, early morning, the evening meal and the morning sun.}

Complete the last lesson of chapter 72: I can say the afternoon, the middle of the night, early morning, the evening meal and the morning sun.
\end{chapteropening}

\section[\texorpdfstring{\ja{ごご} (gogo)}{ごご (gogo)}]{\ja{ごご} (gogo) — the afternoon, p.m.}

\label{lesson:JA-C72-gogo}



\begin{quote}
Before the new one: say the Japanese for a nephew, then the Japanese for a niece.

\emph{Read it:} \textbf{\ja{もういちど、おねがいします}} (Once more, please.) without stopping
\end{quote}

\subsection*{You'll want to know: \ja{ごご}}
\textbf{\ja{ごご}} — \emph{gogo} — \textquotedblleft{}the afternoon, p.m.\textquotedblright{}.

\begin{cousinweb}
\textbf{\ja{ごご}} (\emph{gogo}) is Sino-Japanese (kango): two Chinese characters read with their on'yomi, readings borrowed from Chinese. They are \emph{go}, \textquotedblleft{}noon\textquotedblright{} (Middle Chinese \emph{nguX}) and \emph{go}, \textquotedblleft{}after\textquotedblright{} (\emph{huwX}) (Unihan, KANJIDIC2; Middle Chinese after Baxter–Sagart).
\end{cousinweb}

\begin{grammarlens}[title={What you've built}]
One more word for this chapter.
\end{grammarlens}

\subsection*{Your turn}
\begin{itemize}
\raggedright
  \item \emph{Say it:} \emph{gogo}
  \item \emph{Say it:} \emph{gogo}, once more
  \item \emph{Say it:} say \emph{gogo}
  \item \emph{Recall:} say the Japanese for a nephew, then the Japanese for a niece, then say \emph{gogo} again
\end{itemize}

\begin{tcolorbox}[breakable,colback=teal!4,colframe=teal!35!black,title={Before you move on}]
What does \ja{ごご} mean? (\textquotedblleft{}The afternoon, p.m.\textquotedblright{}.) Say it once more.
\end{tcolorbox}
\section[\texorpdfstring{\ja{よなか} (yonaka)}{よなか (yonaka)}]{\ja{よなか} (yonaka) — the middle of the night}

\label{lesson:JA-C72-yonaka}



\begin{quote}
Before the new one: say the Japanese for a niece, then the Japanese for the afternoon.
\end{quote}

\subsection*{You'll want to know: \ja{よなか}}
\textbf{\ja{よなか}} — \emph{yonaka} — \textquotedblleft{}the middle of the night\textquotedblright{}.

\begin{cousinweb}
\textbf{\ja{よなか}} (\emph{yonaka}) is a native compound, inherited rather than borrowed: \emph{yo}, \textquotedblleft{}night\textquotedblright{}, and \emph{naka}, \textquotedblleft{}middle\textquotedblright{}, both kun'yomi of its kanji (JMdict, KANJIDIC2).
\end{cousinweb}

\begin{grammarlens}[title={What you've built}]
One more word for this chapter.
\end{grammarlens}

\subsection*{Your turn}
\begin{itemize}
\raggedright
  \item \emph{Say it:} \emph{yonaka}
  \item \emph{Say it:} \emph{yonaka}, once more
  \item \emph{Say it:} say \emph{yonaka}
  \item \emph{Recall:} say the Japanese for a niece, then the Japanese for the afternoon, then say \emph{yonaka} again
\end{itemize}

\begin{tcolorbox}[breakable,colback=teal!4,colframe=teal!35!black,title={Before you move on}]
What does \ja{よなか} mean? (\textquotedblleft{}The middle of the night\textquotedblright{}.) Say it once more.
\end{tcolorbox}
\section[\texorpdfstring{\ja{あさがた} (asagata)}{あさがた (asagata)}]{\ja{あさがた} (asagata) — early morning}

\label{lesson:JA-C72-asagata}



\begin{quote}
Before the new one: say the Japanese for the afternoon, then the Japanese for the middle of the night.
\end{quote}

\subsection*{You'll want to know: \ja{あさがた}}
\textbf{\ja{あさがた}} — \emph{asagata} — \textquotedblleft{}early morning\textquotedblright{}.

\begin{cousinweb}
\textbf{\ja{あさがた}} (\emph{asagata}) is a native compound, inherited rather than borrowed: \emph{asa}, \textquotedblleft{}morning\textquotedblright{}, and \emph{kata}, \textquotedblleft{}side; time\textquotedblright{}, voiced to \emph{gata}, both kun'yomi of its kanji (JMdict, KANJIDIC2).
\end{cousinweb}

\begin{grammarlens}[title={What you've built}]
One more word for this chapter.
\end{grammarlens}

\subsection*{Your turn}
\begin{itemize}
\raggedright
  \item \emph{Say it:} \emph{asagata}
  \item \emph{Say it:} \emph{asagata}, once more
  \item \emph{Say it:} say \emph{asagata}
  \item \emph{Recall:} say the Japanese for the afternoon, then the Japanese for the middle of the night, then say \emph{asagata} again
\end{itemize}

\begin{tcolorbox}[breakable,colback=teal!4,colframe=teal!35!black,title={Before you move on}]
What does \ja{あさがた} mean? (\textquotedblleft{}Early morning\textquotedblright{}.) Say it once more.
\end{tcolorbox}
\section[\texorpdfstring{\ja{ゆうはん} (yuuhan)}{ゆうはん (yuuhan)}]{\ja{ゆうはん} (yuuhan) — the evening meal}

\label{lesson:JA-C72-yuuhan}



\begin{quote}
Before the new one: say the Japanese for the middle of the night, then the Japanese for early morning.
\end{quote}

\subsection*{You'll want to know: \ja{ゆうはん}}
\textbf{\ja{ゆうはん}} — \emph{yuuhan} — \textquotedblleft{}the evening meal\textquotedblright{}.

\begin{cousinweb}
\textbf{\ja{ゆうはん}} (\emph{yuuhan}) mixes the two layers: native \emph{yū}, \textquotedblleft{}evening\textquotedblright{}, read as a kun'yomi, and \emph{han}, \textquotedblleft{}cooked rice; a meal\textquotedblright{}, the on'yomi of its second character, borrowed from Chinese (Middle Chinese \emph{bjonX}; JMdict, KANJIDIC2).
\end{cousinweb}

\begin{grammarlens}[title={What you've built}]
One more word for this chapter.
\end{grammarlens}

\subsection*{Your turn}
\begin{itemize}
\raggedright
  \item \emph{Say it:} \emph{yuuhan}
  \item \emph{Say it:} \emph{yuuhan}, once more
  \item \emph{Say it:} say \emph{yuuhan}
  \item \emph{Recall:} say the Japanese for the middle of the night, then the Japanese for early morning, then say \emph{yuuhan} again
\end{itemize}

\begin{tcolorbox}[breakable,colback=teal!4,colframe=teal!35!black,title={Before you move on}]
What does \ja{ゆうはん} mean? (\textquotedblleft{}The evening meal\textquotedblright{}.) Say it once more.
\end{tcolorbox}
\section[\texorpdfstring{\ja{あさひ} (asahi)}{あさひ (asahi)}]{\ja{あさひ} (asahi) — the morning sun}

\label{lesson:JA-C72-asahi}



\begin{quote}
Before the new one: say the Japanese for early morning, then the Japanese for the evening meal.
\end{quote}

\subsection*{You'll want to know: \ja{あさひ}}
\textbf{\ja{あさひ}} — \emph{asahi} — \textquotedblleft{}the morning sun\textquotedblright{}.

\begin{cousinweb}
\textbf{\ja{あさひ}} (\emph{asahi}) is a native compound, inherited rather than borrowed: \emph{asa}, \textquotedblleft{}morning\textquotedblright{}, and \emph{hi}, \textquotedblleft{}sun\textquotedblright{}, both kun'yomi of its kanji (JMdict, KANJIDIC2).
\end{cousinweb}

\begin{grammarlens}[title={What you've built}]
One more word for this chapter.
\end{grammarlens}

\subsection*{Your turn}
\begin{itemize}
\raggedright
  \item \emph{Say it:} \emph{asahi}
  \item \emph{Say it:} \emph{asahi}, once more
  \item \emph{Say it:} say \emph{asahi}
  \item \emph{Recall:} say the Japanese for early morning, then the Japanese for the evening meal, then say \emph{asahi} again
\end{itemize}

\begin{tcolorbox}[breakable,colback=teal!4,colframe=teal!35!black,title={Before you move on}]
What does \ja{あさひ} mean? (\textquotedblleft{}The morning sun\textquotedblright{}.) Say it once more.
\end{tcolorbox}
